% Source: https://artofproblemsolving.com/community/c3h1696278_triangle_and_inequality
% Stable identity: aops-topic-1696278-q01
Write $AD=x\cdot AB$, $AE=y\cdot AC$ and $DF=t\cdot DE$ with $x,y,t\in[0,1]$.

\begin{asy}
usepackage("amsmath");
size(6cm);
defaultpen(fontsize(10pt));
pair A=(1.3,4.2),B=(0,0),C=(6,0);
pair D=A+0.55*(B-A),Ep=A+0.45*(C-A),F=D+0.5*(Ep-D);
fill(D--F--B--cycle,paleblue);fill(F--Ep--C--cycle,palered);
draw(A--B--C--cycle,linewidth(0.9));
draw(D--Ep,linewidth(0.7));draw(F--B,linewidth(0.6));draw(F--C,linewidth(0.6));
dot(D);dot(Ep);dot(F);
label("$A$",A,N);label("$B$",B,SW);label("$C$",C,SE);
label("$D$",D,W);label("$E$",Ep,NE);label("$F$",F,N);
label("\mbox{}",(min(currentpicture,user=true)+max(currentpicture,user=true))/2,black); // keeps the text after the figure black when LaTeX typesets the labels inline
\end{asy}

\textbf{The two areas.} Triangles with a common vertex and bases on one line have areas in the ratio of the bases. So
\[
[DFB]=t\,[DEB]=t(1-x)\,[AEB]=t(1-x)\,y\,S,
\]
since $DF=t\,DE$, $DB=(1-x)AB$ and $AE=y\,AC$. Similarly
\[
[FEC]=(1-t)\,[DEC]=(1-t)(1-y)\,[ADC]=(1-t)(1-y)\,x\,S .
\]

\textbf{Bound.} The minimum of two nonnegative numbers is at most their geometric mean, so
\[
\min\{[DFB],[FEC]\}\le\sqrt{[DFB]\,[FEC]}=S\sqrt{t(1-t)\cdot x(1-x)\cdot y(1-y)}\le S\sqrt{\left(\dfrac{1}{4}\right)^3}=\dfrac {S}{8},
\]
using $u(1-u)\le\dfrac{1}{4}$ (AM-GM) for each of $t,x,y$.

\textbf{Sharpness.} If $D$, $E$, $F$ are the midpoints of $AB$, $AC$, $DE$ ($x=y=t=\dfrac{1}{2}$), both areas equal $\dfrac{1}{2}\cdot\dfrac{1}{2}\cdot\dfrac{1}{2}S=\dfrac {S}{8}$. So the bound cannot be improved, and $k=\boxed{\dfrac{1}{8}}$.
